% CPSY 1291 -- Recitation handout 12
% Compile:  latexmk -lualatex handout-12-final-review.tex
\documentclass[11pt]{article}
\usepackage{cpsy1291-handout}

\begin{document}

\handouthead{12}{Exam 2 review}
  {Lectures 9--19, compressed, and seventeen problems}
  {Week 12 \ $\cdot$\ Exam 2 is in class, Tue 12/1 --- the first day back from recess}

\section*{How to use this}
\addcontentsline{toc}{section}{How to use this}

\hl{Exam 2 is Tuesday 12/1, the first day back from Thanksgiving recess.}
This session is held before recess because there is no session between recess and
the exam. Plan accordingly.

Exam 2 focuses on Lectures 9--19; earlier material appears only as tools. Read
Section~1 first --- it is those eleven lectures compressed into the statements
you should be able to make without notes --- then close it and do Section~2.

\begin{keybox}
\ask{The shape of the second half of this course.} Lectures 1--8 built machinery
for describing a representation. Lectures 9--19 learn representations and keep
asking one question: \emph{does this model tell us anything about the brain, and
how would we know?} Almost every lecture presents a method for answering that,
then presents a reason the method is not enough. Both halves are examinable, and
the second half is where the marks are.
\end{keybox}

\section{Lectures 9--19, compressed}

\subsection{L09 --- Hidden layers and backpropagation}

Every learning rule of Lecture 8 moves the weights against the gradient of its
loss, and the chain rule gives that gradient. One neuron draws one line, so XOR
needs a hidden layer: the hidden units re-represent the four inputs so that one
line separates them. Without a nonlinearity, stacking adds nothing --- a
weighted sum of weighted sums is a weighted sum. With hidden layers the loss is
no longer a bowl, so different random starts can end in different minima.

\textbf{Backpropagation}: (1) a forward pass that keeps every $z$ and $a$;
(2) the output error $\hat y-y$; (3) each hidden unit's error,
$\delta_j=g'(z_j)\sum_k w_{kj}\delta_k$, sent backward layer by layer; (4) every
gradient is error $\times$ input. For many classes, one output per class and a
softmax. The training loop is forward, loss, \texttt{backward()},
\texttt{step()}; autograd runs the backward pass for you. Freezing trained layers
reuses them for a new task.

\subsection{L10 --- Training at scale, generalization and characterizing units}

\textbf{Minibatches} give a noisy, cheap estimate of the gradient;
\textbf{schedules} lower the learning rate during training; \textbf{Adam} keeps
a running average of the gradient (momentum) and divides each weight's step by
the typical size of its own gradient. Sweep the learning rate on a log scale:
too large a step can switch a ReLU unit off on every input, a dead unit.

\textbf{Characterizing a unit}: tuning curves (rectify before averaging over
phase), Gabor fits placed beside V1 cells, object selectivity, and
\textbf{ablation} with a random control to ask whether a unit is needed.

\textbf{Generalization} is performance on data the model never saw. Three
splits: train to fit, validate to choose settings and when to stop, test once
--- with a bootstrap interval, which covers only the draw of test images. Error
decomposes into bias$^2$ + variance + noise, and the best model balances the
first two. Weight decay, data augmentation, dropout and early stopping reduce
overfitting; choose their settings on validation data.

Beyond the test set: \textbf{shortcut learning} --- a cue that predicts the label
in training but not in the world --- fails on new kinds of images (on
cue-conflict images people chose shape in 95.9\% of cases, ResNet-50 in 22.1\%).
Networks fit \textbf{random labels} (Zhang et al.), so the same network memorizes
or generalizes depending on the labels. \textbf{Double descent}: test error peaks
near the interpolation threshold, where a network can first fit its training
set, then falls again. \textbf{Grokking}: validation accuracy can rise long after
training accuracy reaches 100\%. \textbf{Compositional generalization} stays hard
unless training asks for it.

\subsection{L11 --- Learning without labels; the visual system and CNNs}

An \textbf{autoencoder} learns a code by reconstructing its input, so it needs no
labels. A linear autoencoder with tied weights on centered data minimizes the PCA
objective, so it finds the subspace PCA finds; a nonlinear one can follow a
curved manifold. \textbf{Sparse coding} describes each whitened natural-image
patch with a few atoms from a dictionary; the atoms come out localized, oriented
and bandpass, like V1 simple cells, though nothing about edges was built in.
Measured on Ringach's plot, their width matches V1 and their length does not.

The ventral stream: center--surround, then oriented \textbf{simple cells} (an
oriented template at one position: selectivity), then \textbf{complex cells}
(the same orientation anywhere in a larger field, like a maximum over simple
cells: invariance), repeated up to IT, where a linear readout reports the object.
HMAX alternates the two with hand-set filters. A CNN learns them: filters shared
across positions (convolution), stride and pooling to shrink the maps, global
average pooling to a feature vector, a readout; its first layer learns V1-like
filters. Residual blocks and BatchNorm make deep networks trainable (training
mode uses batch statistics, evaluation mode the stored ones). Reuse: freeze the
encoder or fine-tune its last block. Check what a network relies on: texture
bias, occlusion.

\subsection{L12 --- Comparing networks with the brain}

\textbf{Linear predictivity}: a ridge readout from a network stage to each
recorded site, scored on held-out images and divided by the site's split-half
reliability, so 1 means the model predicts the site as well as the site predicts
itself; report the median over sites. Deeper stages predict later areas (V4
peaks one stage before IT). \textbf{RSA} fits nothing and compares geometry;
predictivity fits a weight per unit, so reweighting units barely changes it.
Report both, each against its own ceiling.

Networks that recognize better predicted IT better (Yamins et al.\ 2014), but
past a point accuracy stops helping (Brain-Score), many very different networks
reach the same score (Conwell et al.\ 2024), and the scores disagree with each
other. Stronger tests: \textbf{reverse predictivity} (networks predict about half
of IT, IT predicts only 5--20\% of a network's units) and \textbf{control} with
synthesized images. Beyond neurons: the image regions people use (ClickMe), and
\textbf{decoding} --- a linear readout pointed at a brain (Kamitani \& Tong).

\subsection{L13 --- Understanding what networks compute}

\textbf{Attribution maps} --- occlusion, saliency (the gradient at the pixels),
Grad-CAM (the gradient at the last feature maps) --- say which pixels moved one
class's score, and each answers a slightly different question. A map is a
hypothesis. \textbf{Interventions} test the map: inpaint the region and compare
with the same shape at random places. \textbf{Sanity checks} test the method:
randomize the weights from the top down; gradient and Grad-CAM change and pass,
guided backpropagation does not and fails. A map shows where, not what: the
``elephant'' map covered the cat, and only swapping the texture showed that the
skin decided. Shortcuts can be planted and caught: a map suggests the cue (snow,
a hospital marker), an intervention confirms it. Map, then intervention, then
control.

\subsection{L14 --- Generative models: VAEs, GANs and diffusion}

A generative model learns $p(\mathbf{x})$ and samples new data, usually by
drawing a latent $\mathbf{z}$ from a simple prior and mapping it to $\mathbf{x}$.
\textbf{VAE}: maximize the ELBO, reconstruction minus the KL divergence of each
input's code distribution from the prior; the KL is what fills the holes a plain
autoencoder leaves. The \textbf{reparameterization trick},
$\mathbf{z}=\boldsymbol{\mu}+\boldsymbol{\sigma}\odot\boldsymbol{\epsilon}$, moves
the randomness outside the network so gradients reach the encoder.
\textbf{GAN}: a generator minimizes and a discriminator maximizes one objective;
at equilibrium $p_G=p_{\text{data}}$ and $D=\tfrac12$ everywhere. Samples are
sharp, training unstable, and the generator can collapse onto a few outputs;
steerable latents make tools such as GANalyze. \textbf{Diffusion}: a fixed forward
process adds noise, with a closed form that jumps to any step; a network learns
to predict the added noise with a plain squared error; sampling takes many small
denoising steps, which is why it is slow.

\subsection{L15 --- RNNs, sequences and decisions}

An RNN carries a state,
$\mathbf{h}^{\langle t\rangle}=g(\mathbf{W}_{hh}\mathbf{h}^{\langle t-1\rangle}+\mathbf{W}_{hx}\mathbf{x}^{\langle t\rangle}+\mathbf{b}_h)$,
with the same weights at every step; unrolled, it is a $T$-deep network, trained
with a loss summed over time and backpropagation through time. Units leak at the
rate $\Delta t/\tau$.

Decisions: \textbf{psychometric} and \textbf{chronometric} functions, the Weibull
threshold, evidence accumulated to a \textbf{bound}, and the speed--accuracy
trade-off. LIP neurons ramp toward a common level. A task-trained RNN's log odds
can play the accumulated evidence, so its choices, reaction times and units can
be compared with a monkey's.

Long memories are hard because the gradient travels back through a
\emph{product} of Jacobians, $\mathrm{diag}(g')\,\mathbf{W}_{hh}$, one per step.
Spectral radius below 1: it vanishes; above 1: it can explode. Clipping fixes
exploding but cannot create gradient; \textbf{gating} (LSTM, GRU) fixes vanishing,
because the cell is updated by addition.

\subsection{L16 --- Neural dynamics and attractors}

With its input held fixed, a recurrent network is a dynamical system that rests
at \textbf{fixed points}. Point attractors store memories and choices; a
\textbf{line attractor} --- a line of fixed points, with no restoring force along
it --- holds and sums a value: working memory, evidence accumulation, eye
position. A decision is a fall into one basin. To analyse a trained network,
minimize the speed $q(\mathbf{h})=\tfrac12\|F(\mathbf{h})-\mathbf{h}\|^2$ to find
fixed points, then classify each by the eigenvalues of its Jacobian: all
$|\lambda|<1$ stable, some above 1 a saddle, $|\lambda|$ near 1 a slow direction.

\textbf{Hopfield networks} store patterns in Hebbian weights and recall by
descending an energy. Stored patterns interfere, and mixtures become minima
nobody stored; random patterns are recalled up to about $0.138\,D$, and correlated
patterns fail far earlier. A \textbf{modern Hopfield} network recalls with a
softmax, stores very many memories, and its update is one step of attention.

\subsection{L17 --- Attention and the transformer}

Each position's new vector is a weighted average of value rows,
$\mathrm{softmax}(\mathbf{Q}\mathbf{K}^{\!\top}/\sqrt{d_k})\,\mathbf{V}$, with
$\mathbf{Q},\mathbf{K},\mathbf{V}$ projections of the same input. The
$\sqrt{d_k}$ keeps large scores from saturating the softmax. Any position reaches
any other in one step, at a cost of $T\times T$ scores. Multi-head attention runs
several heads side by side; positional embeddings restore order; a block is
attention and a per-position MLP, each with a residual and a LayerNorm. A
language model predicts the next symbol at every position at once, the causal
mask keeps each prediction from seeing its answer, and a bigram table is the
baseline. One modern-Hopfield update is one attention step. An attention map
shows where a head reads, not why the output is what it is.

\subsection{L18 --- LLMs, self-supervised learning and predictive coding}

One objective (next token), one architecture (a decoder transformer) and scale
make an LLM. \textbf{Scaling laws}: loss is a predictable power law in
parameters and data; compute-optimal training scales data with parameters
(Chinchilla, about 20 tokens per parameter). \textbf{In-context learning}
appears without being trained for; \textbf{emergence} is partly a property of the
metric --- all-or-nothing scores make gradual improvement look like a jump.
\textbf{Surprisal} from LLMs tracks brain responses to language, but very large
models fit reading times worse.

\textbf{Self-supervised learning} makes the data label itself:
\textbf{contrastive} learning pulls two augmented views of one image together and
pushes other images away; \textbf{masked prediction} hides part of the input and
reconstructs it. Self-supervised vision models rival supervised ones at
predicting the ventral stream. \textbf{Predictive coding}: the brain may predict
its input and send errors up (Rao \& Ballard).

\subsection{L19 --- World models and frontiers}

A \textbf{world model} is a predictive model of an environment, for planning and
imagination. Whether prediction alone builds one is open: a transformer trained
on taxi routes predicts valid turns almost perfectly while its implied street map
is wrong. Current AI breaks on robustness and shortcuts, on visual and abstract
reasoning, and on data efficiency. And the closing question: \emph{does better AI
mean better models of biology?} Alignment with brains and with human strategies
can fall as accuracy rises. The proposed fixes point the same way: rethink the
data diet, align training to human strategy (ClickMe, harmonization), and bring
back recurrence and feedback.

\section{Problems}

\ask{1.} A network reaches 100\% training accuracy on randomly shuffled labels.
Why does this undermine capacity-based explanations of generalization?

\ask{2.} Sketch a double-descent curve. Mark the interpolation threshold and say
what is true of the network there.

\ask{3.} A classifier reaches 94\% on a held-out test set and 41\% on X-rays from
a different hospital. Name the phenomenon and say how the tools of L13 would
catch it.

\ask{4.} Why does a linear autoencoder with tied weights find the PCA subspace,
and why is ``the PCA subspace'' the right phrasing rather than ``the principal
components''?

\ask{5.} Sparse coding on whitened natural images produces Gabor-like filters.
Explain in two sentences why this is evidence about V1 rather than a coincidence,
and say where the match fails.

\ask{6.} What does max-pooling buy, and which cortical cell type is it likened
to?

\ask{7.} A network stage predicts an IT site with held-out $R^2=0.40$; the site's
split-half reliability is 0.80. Give the score, say what 1 would mean, and say
why RSA should be reported as well.

\ask{8.} Networks predict about half of IT's repeatable variance, but IT
predicts only 5--20\% of a network's units. What does this reverse test show?

\ask{9.} You randomize a network's weights from the top layer down and its
attribution map barely changes. What follows?

\ask{10.} Name the two terms of the ELBO and what each does, and say why the
reparameterization trick is needed.

\ask{11.} An RNN trains on sequences of length 20 and fails at length 300.
Name the mechanism, and say why gradient clipping does not fix it.

\ask{12.} A trained RNN has a fixed point with one Jacobian eigenvalue at 1.00
and the rest near 0.3. What computation does this support, and give a task that
would need it.

\ask{13.} A Hopfield network of 400 units is asked to store 100 random patterns.
What happens on recall, and what would change with correlated patterns such as
digits?

\ask{14.} Why is the $\sqrt{d_k}$ in scaled dot-product attention necessary, and
what fails without it?

\ask{15.} A colleague shows an attention map and claims it explains the model's
answer. Give two objections and the test that would settle it.

\ask{16.} A task sits near chance for small language models and then jumps past
a scale threshold. Give the alternative explanation from L18.

\ask{17.} State the closing question of L19 and the evidence that makes it
worrying.

\subsection*{Answers}

\ask{1.} If a network can fit arbitrary labels, its capacity is sufficient to
memorize the training set outright --- so capacity cannot be what prevents it from
memorizing when the labels are real. The same network memorizes or generalizes
depending on the labels, so whatever explains generalization must involve the
data, the optimizer, or the architecture, not the parameter count alone.

\ask{2.} Test error falls, rises to a peak near the interpolation threshold, then
falls again as the network grows. At the threshold the network can just fit its
training set. Larger networks can fit it in many ways, and which fit training
finds affects the test error.

\ask{3.} Shortcut learning: the model used a cue that predicts the label in the
training images but not in the world, such as a hospital's marker in the corner.
An attribution map suggests the cue; an intervention confirms it (edit the region
and watch the decision), with a control (the same edit to a same-size region
elsewhere).

\ask{4.} With tied weights, centered data and no nonlinearity, the reconstruction
is $\mathbf{W}^\top\mathbf{W}\mathbf{x}$, which lies in the $k$-dimensional
subspace spanned by the rows of $\mathbf{W}$; minimizing the squared
reconstruction error is then the PCA objective. It finds the \emph{subspace}
rather than the components because any rotation within that subspace
reconstructs equally well --- nothing in the objective orders or aligns the axes.

\ask{5.} The filters were not fitted to neural data, and nothing about edges was
built in: they fall out of describing natural images with a few active atoms,
and come out localized, oriented and bandpass like V1 simple cells. Measured on
Ringach's plot, the width matches (0.33 against 0.32) but the length does not
(0.87 against 0.41): the objective predicts V1 to first order, and measuring shows
where it falls short.

\ask{6.} It keeps the largest value in each small window, which tolerates small
shifts and shrinks the map. It is likened to the complex cell, which behaves like
the maximum over simple cells of the same orientation at nearby positions.

\ask{7.} $0.40/0.80=0.50$. A score of 1 means the model predicts the site as well
as the site predicts itself. RSA assumes something different --- it fits no
weights, so it is sensitive to how units are weighted --- and the two can pick
different stages, so report both, each against its own ceiling.

\ask{8.} Good networks carry many dimensions that IT lacks. A model can predict a
brain area well and still be a much larger representation than the area; a
forward score alone cannot see that. (Monkey to monkey, the fit is about 40\%
each way.)

\ask{9.} The map is not explaining what the network learned; it is driven by the
image and the architecture. In the lecture, gradient and Grad-CAM maps change and
pass, while guided backpropagation still shows the bird and fails. Any claim such
a map supported about the trained model is unsupported.

\ask{10.} Reconstruction: codes drawn from the encoder should decode back to
$\mathbf{x}$. KL: each input's code distribution should stay close to the prior,
which fills the holes so that samples decode sensibly. Backprop cannot pass
through a random draw; writing
$\mathbf{z}=\boldsymbol{\mu}+\boldsymbol{\sigma}\odot\boldsymbol{\epsilon}$ moves
the randomness into $\boldsymbol{\epsilon}$, so $\mathbf{z}$ is a differentiable
function of $\boldsymbol{\mu}$ and $\boldsymbol{\sigma}$.

\ask{11.} Vanishing gradients: the error travels back through a product of 300
Jacobians, and with spectral radius below 1 the product shrinks. Clipping rescales
gradients that are too \emph{large} and cannot create gradient; the fix is gating,
which updates the cell by addition.

\ask{12.} A line attractor --- a slow direction that neither decays nor grows, with
everything else contracting onto it. It holds and sums a value: accumulating
evidence for a decision, or holding a graded value in working memory, such as eye
position.

\ask{13.} Random patterns are recalled up to about $0.138\times400\approx55$; at
100, the stored patterns stop being minima and recall collapses into mixtures.
Correlated patterns, such as faces, words or digits, fail far earlier; this
limit motivated the modern Hopfield network's sharper energy.

\ask{14.} A dot product of two random $d_k$-dimensional vectors has a spread that
grows like $\sqrt{d_k}$, so without the scaling the softmax argument becomes
large, the softmax puts almost all weight on one position, and its gradient
becomes close to zero. Training then fails for large $d_k$.

\ask{15.} (i) There are many heads and many layers, so ``the'' map is one chosen
among many. (ii) A map shows where a head reads, not why the output is what it
is. The test is an intervention: change the input or switch the head off and see
whether the answer changes.

\ask{16.} The metric. Most such jumps come from all-or-nothing scores such as
exact match; scored with a smooth metric, the same models improve gradually
(Schaeffer et al.).

\ask{17.} Does better AI mean better models of biology? Past a point, more
accurate networks do not predict IT better, and alignment with brains and with
human strategies can fall as accuracy rises. The proposed responses: compare
strategies rather than outputs, align training to human strategy, rethink the
data diet, and bring back recurrence and feedback.

\section*{The night before}
\addcontentsline{toc}{section}{The night before}

Eight things worth being able to state cold:

\begin{enumerate}
  \item the four steps of backpropagation, and why linear layers collapse into
        one;
  \item the three splits, and what the random-labels result rules out;
  \item what a tied linear autoencoder finds, and what sparse coding of natural
        images finds;
  \item how a predictivity score is normalized, and why one score is not enough;
  \item map, intervention, control --- and what a sanity check tests;
  \item the ELBO's two terms, and what a diffusion network predicts;
  \item vanishing gradients as a product, why gating and not clipping fixes it,
        and what a line attractor is;
  \item the attention equation, the causal mask, and what an attention map does
        not show.
\end{enumerate}

\end{document}
